\documentclass[12pt]{article}
\usepackage[spanish]{babel}
\usepackage[utf8]{inputenc}
\usepackage[T1]{fontenc}
\usepackage{geometry}
\usepackage{graphicx}
\usepackage{setspace}
\usepackage{titlesec}
\usepackage{enumitem}

\geometry{margin=1in}
\onehalfspacing

\begin{document}

\begin{titlepage}
    \centering

    % Logo opcional
    %\includegraphics[width=0.3\textwidth]{logo.png}\par\vspace{1cm}

    {\scshape\Large Redes \par}
    \vspace{1.5cm}

    {\huge\bfseries Cuestionario QoS \par}
    \vspace{0.5cm}

    \vfill

    {\large
    Nombre del Alumno: Emilio Izquierdo Montero {\hspace{6cm}} \\
    \vspace{0.3cm}
    Profesor: Pedro Pablo Martinez Carmona {\hspace{6cm}} \\
    \vspace{0.3cm}
    Grupo / Semestre: 8 {\hspace{6cm}} \\
    \vspace{0.3cm}
    Fecha: 6 de mayo de 2026 {\hspace{6cm}}
    \par}

    \vfill

\end{titlepage}

\section*{INTRODUCCIÓN A QoS - CALIDAD DE SERVICIO}

La Calidad de Servicio (QoS, por sus siglas en inglés \textit{Quality of Service}) es un conjunto de tecnologías y mecanismos utilizados en redes de computadoras para administrar y priorizar el tráfico de datos. Su objetivo principal es garantizar que ciertos servicios o aplicaciones reciban un mejor rendimiento dentro de la red, especialmente en situaciones donde existe congestión o limitaciones de ancho de banda.

QoS es ampliamente utilizado en redes empresariales, servidores, sistemas de telecomunicaciones y administración de infraestructura de TI. Gracias a esta tecnología es posible priorizar aplicaciones críticas como videollamadas, VoIP, transmisión de video, servicios en tiempo real y aplicaciones sensibles a la latencia.

Entre los parámetros más importantes que QoS controla se encuentran:

\begin{itemize}
    \item \textbf{Ancho de banda:} cantidad de datos que pueden transmitirse por la red.
    \item \textbf{Latencia:} tiempo que tarda un paquete en llegar a su destino.
    \item \textbf{Jitter:} variación en el tiempo de llegada de los paquetes.
    \item \textbf{Pérdida de paquetes:} paquetes que no llegan correctamente al destino.
\end{itemize}

En administración de sistemas y redes, QoS permite optimizar recursos y mantener estabilidad en los servicios críticos. Por ejemplo, en una empresa se puede configurar QoS para que el tráfico de videoconferencias tenga mayor prioridad que las descargas de archivos o navegación web.

También es común encontrar QoS en routers, switches administrables y sistemas operativos modernos, donde se aplican reglas para clasificar y controlar el tráfico de red según políticas definidas por el administrador.

\section*{Preguntas}

\subsection*{1. Pregunta abierta}
\textbf{¿Qué significa QoS?}

Respuesta: Quality of Service o Calidad de Servicio.

\vspace{0.5cm}

\subsection*{2. Falso o Verdadero}
\textbf{QoS ayuda a priorizar tráfico importante dentro de una red.}

Respuesta: Verdadero.

\vspace{0.5cm}

\subsection*{3. Opción múltiple}
\textbf{¿Cuál de los siguientes servicios requiere QoS?}

\begin{enumerate}[label=\alph*)]
    \item Calculadora
    \item Videollamada
    \item Bloc de notas
    \item Explorador de archivos
\end{enumerate}

Respuesta: b) Videollamada.

\vspace{0.5cm}

\subsection*{4. Relaciona columnas}

\begin{tabular}{|l|l|}
\hline
Latencia & Tiempo de llegada \\ \hline
Jitter & Variación de retraso \\ \hline
Ancho de banda & Capacidad de transmisión \\ \hline
\end{tabular}

\vspace{0.5cm}

\subsection*{5. Completa la oración}
\textbf{QoS se utiliza para mejorar el rendimiento de la \_\_\_\_\_\_.}

Respuesta: red.

\vspace{0.5cm}

\subsection*{6. Pregunta abierta}
\textbf{¿Qué dispositivo de red puede implementar QoS?}

Respuesta: Routers y switches administrables.

\vspace{0.5cm}

\subsection*{7. Falso o Verdadero}
\textbf{QoS únicamente funciona en redes inalámbricas.}

Respuesta: Falso.

\vspace{0.5cm}

\subsection*{8. Opción múltiple}
\textbf{¿Qué parámetro mide la pérdida de paquetes?}

\begin{enumerate}[label=\alph*)]
    \item Packet Loss
    \item CPU Usage
    \item RAM
    \item DNS
\end{enumerate}

Respuesta: a) Packet Loss.

\vspace{0.5cm}

\subsection*{9. Subrayar}
\textbf{Subraya el servicio más importante para QoS:}

\begin{itemize}
    \item Juegos sin conexión
    \item Voz sobre IP (VoIP)
    \item Editor de texto
\end{itemize}

Respuesta: Voz sobre IP (VoIP).

\vspace{0.5cm}

\subsection*{10. Pregunta abierta}
\textbf{¿Por qué es importante QoS en redes empresariales?}

Respuesta: Porque permite garantizar estabilidad y prioridad a servicios críticos.

\end{document}
